\documentclass[12pt,a4paper]{article}

% --- Paquetes ---
\usepackage[utf8]{inputenc}
\usepackage[T1]{fontenc}
\usepackage[spanish]{babel}
\usepackage{geometry}
\usepackage{graphicx}
\usepackage{hyperref}
\usepackage{listings}
\usepackage{xcolor}
\usepackage{booktabs}
\usepackage{enumitem}
\usepackage{fancyhdr}
\usepackage{titlesec}
\usepackage{parskip}

\geometry{margin=2.5cm}

% --- Colores ---
\definecolor{codebg}{RGB}{245,245,245}
\definecolor{codeframe}{RGB}{180,180,180}
\definecolor{codecomment}{RGB}{0,128,0}
\definecolor{codestring}{RGB}{163,21,21}
\definecolor{codekeyword}{RGB}{0,0,200}
\definecolor{linkcolor}{RGB}{0,80,160}

% --- Hyperref ---
\hypersetup{
    colorlinks=true,
    linkcolor=linkcolor,
    urlcolor=linkcolor,
    citecolor=linkcolor,
}

% --- Listings ---
\lstdefinestyle{bashstyle}{
    language=bash,
    backgroundcolor=\color{codebg},
    frame=single,
    rulecolor=\color{codeframe},
    basicstyle=\ttfamily\small,
    keywordstyle=\color{codekeyword}\bfseries,
    commentstyle=\color{codecomment}\itshape,
    stringstyle=\color{codestring},
    breaklines=true,
    breakatwhitespace=false,
    showstringspaces=false,
    tabsize=2,
    xleftmargin=1em,
    framexleftmargin=0.5em,
    aboveskip=1em,
    belowskip=1em,
    literate={á}{{\'a}}1 {é}{{\'e}}1 {í}{{\'i}}1 {ó}{{\'o}}1 {ú}{{\'u}}1
             {ñ}{{\~n}}1 {Ñ}{{\~N}}1,
    morekeywords={sudo,docker,sysctl,curl,chmod,mkdir,systemctl,usermod,git,pip3,cd,bash,python3},
}

\lstdefinestyle{envstyle}{
    backgroundcolor=\color{codebg},
    frame=single,
    rulecolor=\color{codeframe},
    basicstyle=\ttfamily\small,
    breaklines=true,
    showstringspaces=false,
    xleftmargin=1em,
    framexleftmargin=0.5em,
    aboveskip=1em,
    belowskip=1em,
}

\lstset{style=bashstyle}

% --- Header/Footer ---
\pagestyle{fancy}
\fancyhf{}
\fancyhead[L]{\small CTF Layer 2 Security}
\fancyhead[R]{\small Redes de Computadoras}
\fancyfoot[C]{\thepage}
\renewcommand{\headrulewidth}{0.4pt}

% ======================================================================
\begin{document}

% --- Portada ---
\begin{titlepage}
    \centering
    \vspace*{2cm}

    {\LARGE\textbf{Universidad San Francisco de Quito}}\\[0.5cm]
    {\large Colegio de Ciencias e Ingenier\'ias}\\[0.3cm]
    {\large Redes de Computadoras}\\[2.5cm]

    {\Huge\textbf{CTF Layer 2 Security}}\\[0.5cm]
    {\LARGE Gu\'ia de Ejecuci\'on Paso a Paso}\\[3cm]

    {\Large\textbf{Equipo:}}\\[0.5cm]
    {\large
        Julian Leon\\
        James Soto\\
        Mauricio Mantilla\\
    }\\[2.5cm]

    {\large Marzo 2026}

    \vfill
    {\small Repositorio: \url{https://github.com/jjjulianleon/ProyectoFinalRedes}}
\end{titlepage}

% --- Tabla de Contenidos ---
\tableofcontents
\newpage

% ======================================================================
\section{Introducci\'on}

Este documento describe c\'omo clonar, configurar y ejecutar el proyecto \textbf{CTF Layer 2 Security} en una m\'aquina local. El entorno completo corre sobre Docker y simula una red con contenedores para practicar ataques y defensas de Capa~2 (ARP Spoofing, MAC Flooding) con monitoreo mediante Suricata y Wazuh.

% ======================================================================
\section{Prerrequisitos}

\subsection{M\'aquina Virtual con Kali Linux}

Se recomienda una VM de Kali Linux Rolling (la versi\'on m\'as reciente). Opciones de hipervisor:

\begin{itemize}
    \item \textbf{Linux:} virt-manager (recomendado)
    \item \textbf{Windows/Mac:} VirtualBox o VMware
\end{itemize}

Requisitos m\'inimos de la VM:
\begin{itemize}
    \item \textbf{RAM:} 8\,GB asignados a la VM
    \item \textbf{Disco:} al menos 30\,GB
    \item \textbf{Sistema:} Kali Rolling (latest)
\end{itemize}

\subsection{Docker + Docker Compose}

Versiones probadas: Docker CE v27.5.1 y Docker Compose plugin v5.1.0.

Instalaci\'on en Kali:

\begin{lstlisting}
sudo apt update
sudo apt install docker.io
sudo systemctl enable --now docker
sudo usermod -aG docker $USER
\end{lstlisting}

El plugin de Docker Compose no est\'a en los repositorios de Kali, por lo que se debe instalar manualmente:

\begin{lstlisting}
DOCKER_CONFIG=${DOCKER_CONFIG:-$HOME/.docker}
mkdir -p $DOCKER_CONFIG/cli-plugins
curl -SL https://github.com/docker/compose/releases/latest/download/docker-compose-linux-x86_64 \
  -o $DOCKER_CONFIG/cli-plugins/docker-compose
chmod +x $DOCKER_CONFIG/cli-plugins/docker-compose
\end{lstlisting}

Verificar la instalaci\'on:

\begin{lstlisting}
docker compose version
\end{lstlisting}

\textbf{Importante:} despu\'es de ejecutar \texttt{usermod}, cerrar sesi\'on y volver a iniciarla para que el grupo \texttt{docker} surta efecto. Si no, usar \texttt{sudo} antes de cada comando \texttt{docker}.

\subsection{Git}

\begin{lstlisting}
sudo apt install git
\end{lstlisting}

% ======================================================================
\section{Pasos de Ejecuci\'on}

\subsection{Paso 1: Clonar el repositorio}

\begin{lstlisting}
git clone https://github.com/jjjulianleon/ProyectoFinalRedes.git
cd ProyectoFinalRedes/ctf-layer2-security
\end{lstlisting}

\subsection{Paso 2: Crear el archivo \texttt{.env}}

Crear un archivo llamado \texttt{.env} dentro del directorio \texttt{ctf-layer2-security/} con el siguiente contenido:

\begin{lstlisting}[style=envstyle]
FLAG1=FLAG{arp_spoof_mitm_captured}
FLAG2=FLAG{mac_flood_broadcast_leak}
FLAG3=FLAG{plaintext_credentials_exposed}
\end{lstlisting}

Se puede crear con un editor de texto o directamente desde la terminal:

\begin{lstlisting}
cat > .env << 'EOF'
FLAG1=FLAG{arp_spoof_mitm_captured}
FLAG2=FLAG{mac_flood_broadcast_leak}
FLAG3=FLAG{plaintext_credentials_exposed}
EOF
\end{lstlisting}

\subsection{Paso 3: Generar certificados SSL de Wazuh}

\begin{lstlisting}
cd containers/wazuh
sudo docker compose -f generate-certs-compose.yml run --rm generator
cd ../..
\end{lstlisting}

\subsection{Paso 4: Configurar par\'ametros del sistema e iniciar contenedores}

\begin{lstlisting}
sudo sysctl -w vm.max_map_count=262144
sudo docker compose up -d --build
\end{lstlisting}

Esto construye e inicia los \textbf{13 contenedores}. La primera vez tarda aproximadamente \textbf{10 minutos} en descargar las im\'agenes.

\subsection{Paso 5: Verificar el entorno}

\begin{lstlisting}
sudo bash scripts/verify_environment.sh
\end{lstlisting}

Las \textbf{31 verificaciones} deben pasar correctamente.

\subsection{Paso 6: Configurar CTFd}

\begin{enumerate}
    \item Abrir \url{http://localhost:8000} en el navegador.
    \item Completar el asistente de configuraci\'on:
    \begin{itemize}
        \item \textbf{Team name:} admin
        \item \textbf{Email:} admin@ctf.local
        \item \textbf{Password:} (elegir uno)
        \item \textbf{Modo:} Team Mode
    \end{itemize}
    \item Generar un token de API: \texttt{Settings} $\rightarrow$ \texttt{Access Tokens} $\rightarrow$ \texttt{Generate}.
    \item Crear los retos (challenges):
\end{enumerate}

\begin{lstlisting}
# Instalar requests en victim1 y copiar script de setup
sudo docker exec victim1 pip3 install requests 2>/dev/null
sudo docker cp containers/ctfd/config/setup_ctfd.py victim1:/tmp/
sudo docker exec victim1 python3 /tmp/setup_ctfd.py <TU_API_TOKEN>
\end{lstlisting}

Reemplazar \texttt{<TU\_API\_TOKEN>} con el token generado en el paso anterior.

\subsection{Paso 7: Configurar integraci\'on Wazuh + Suricata}

\begin{lstlisting}
sudo bash scripts/setup_wazuh_suricata.sh
\end{lstlisting}

\subsection{Paso 8: Ejecutar la demo del CTF}

\begin{lstlisting}
sudo bash scripts/run_ctf_demo.sh
\end{lstlisting}

Este script tarda aproximadamente \textbf{2 minutos} y realiza lo siguiente:

\begin{enumerate}
    \item Inicia el monitoreo del Blue Team.
    \item Ejecuta ARP Spoofing sobre las 3 v\'ictimas.
    \item Captura las 3 flags.
    \item Ejecuta MAC Flooding.
    \item Restaura las tablas ARP.
    \item Recopila toda la evidencia en el directorio \texttt{evidence\_YYYYMMDD\_HHMMSS/}.
\end{enumerate}

\subsection{Paso 9: Ver resultados}

\subsubsection{Directorio de evidencia (\texttt{evidence\_*/})}

\begin{itemize}
    \item \texttt{arp\_alerts.log} -- Detecciones de ARP spoofing
    \item \texttt{mac\_flood\_alerts.log} -- Detecciones de MAC flooding
    \item \texttt{suricata\_fast.log} -- Alertas del NIDS Suricata
    \item \texttt{captura\_ctf.pcap} -- Captura completa de tr\'afico (abrir con Wireshark)
    \item \texttt{victim*\_response.html} / \texttt{victim*\_credentials.txt} -- Datos capturados
\end{itemize}

\subsubsection{Interfaces web}

\begin{itemize}
    \item \textbf{CTFd:} \url{http://localhost:8000}
    \item \textbf{Wazuh Dashboard:} \url{https://localhost:5601}
    \begin{itemize}
        \item Usuario: \texttt{admin}
        \item Contrase\~na: \texttt{SecretPassword}
    \end{itemize}
\end{itemize}

% ======================================================================
\section{Ejecuci\'on Manual (Modo Interactivo)}

Para una ejecuci\'on m\'as detallada y controlada, se pueden ejecutar los comandos manualmente en terminales separadas.

\subsection{Red Team (ataque)}

\begin{lstlisting}
# Terminal 1: ARP Spoofing
sudo docker exec -it redteam \
  python3 /tools/arp_spoof.py -t 172.20.0.12 -g 172.20.0.2 --broadcast

# Terminal 2: Capturar flags
sudo docker exec -it redteam python3 /tools/capture_flags.py

# Terminal 3: MAC Flooding
sudo docker exec -it redteam python3 /tools/mac_flood.py -c 5000

# Enviar flags al CTFd
sudo docker exec -it redteam \
  python3 /tools/submit_flag.py --list --token <TOKEN>
sudo docker exec -it redteam \
  python3 /tools/submit_flag.py -f "FLAG{...}" -c 1 --token <TOKEN>
\end{lstlisting}

\subsection{Blue Team (defensa)}

\begin{lstlisting}
# Terminal 1: Monitor ARP
sudo docker exec -it blueteam \
  python3 /tools/arp_monitor.py --known 172.20.0.2 172.20.0.10 172.20.0.11 172.20.0.12

# Terminal 2: Detector de anomalias MAC
sudo docker exec -it blueteam python3 /tools/mac_anomaly_detector.py

# Terminal 3: Restaurar ARP (despues del ataque)
sudo docker exec -it blueteam python3 /tools/arp_restore.py --continuous
\end{lstlisting}

% ======================================================================
\section{Arquitectura de Red}

La siguiente tabla muestra los contenedores, sus direcciones IP y roles dentro del entorno:

\begin{table}[h!]
\centering
\begin{tabular}{@{}llp{7cm}@{}}
\toprule
\textbf{Contenedor} & \textbf{IP} & \textbf{Rol} \\
\midrule
Gateway         & 172.20.0.2   & Router / objetivo de ataques \\
Victim 1        & 172.20.0.10  & Servidor HTTP (flag en HTML) \\
Victim 2        & 172.20.0.11  & Servidor de archivos (flag en archivo) \\
Victim 3        & 172.20.0.12  & Agente de monitoreo (flag en tr\'afico) \\
Red Team        & 172.20.0.100 & Atacante \\
Blue Team       & 172.20.0.200 & Defensor \\
Suricata        & (red gateway) & NIDS (Sistema de Detecci\'on de Intrusos) \\
Wazuh Manager   & 172.20.0.240 & SIEM -- Gestor \\
Wazuh Indexer   & 172.20.0.241 & SIEM -- Almacenamiento de logs \\
Wazuh Dashboard & 172.20.0.242 & SIEM -- Interfaz web \\
CTFd            & 172.20.0.250 & Plataforma CTF \\
CTFd DB         & 172.20.0.251 & Base de datos (MariaDB) \\
CTFd Cache      & 172.20.0.252 & Cach\'e (Redis) \\
\bottomrule
\end{tabular}
\caption{Contenedores y direcciones IP del entorno CTF}
\end{table}

% ======================================================================
\section{Soluci\'on de Problemas}

\begin{description}[style=nextline, leftmargin=0pt]

\item[\texttt{permission denied} al usar Docker]
Ejecutar \texttt{sudo usermod -aG docker \$USER} y luego cerrar sesi\'on y volver a iniciarla. Alternativamente, usar \texttt{sudo} antes de cada comando \texttt{docker}.

\item[Error de \texttt{vm.max\_map\_count} (Wazuh)]
Ejecutar antes de iniciar los contenedores:
\begin{lstlisting}
sudo sysctl -w vm.max_map_count=262144
\end{lstlisting}
Este par\'ametro se pierde al reiniciar la m\'aquina; hay que ejecutarlo de nuevo despu\'es de cada reboot.

\item[Los contenedores no inician]
Verificar la RAM disponible con \texttt{free -h}. Se necesitan al menos \textbf{6\,GB libres}.

\item[CTFd no es accesible]
Esperar 30 segundos despu\'es del arranque. Revisar logs con:
\begin{lstlisting}
sudo docker logs ctfd
\end{lstlisting}

\item[Wazuh Dashboard muestra error 503]
Esperar 2--3 minutos para que el indexer se inicialice. Revisar logs con:
\begin{lstlisting}
sudo docker logs wazuh-indexer
\end{lstlisting}

\item[\texttt{docker-compose} no encontrado]
En Kali, el comando es \texttt{docker compose} (con espacio, no gui\'on). Instalar el plugin manualmente como se indica en la secci\'on de prerrequisitos.

\end{description}

% ======================================================================
\section{Comandos \'Utiles}

\begin{lstlisting}
# Ver estado de todos los contenedores
sudo docker ps --format "table {{.Names}}\t{{.Status}}"

# Reconstruir todo desde cero
sudo docker compose down -v && sudo docker compose up -d --build

# Ver logs de un contenedor
sudo docker logs <nombre_contenedor> 2>&1 | tail -20

# Entrar a un contenedor
sudo docker exec -it redteam bash
sudo docker exec -it blueteam bash

# Detener todo
sudo docker compose down
\end{lstlisting}

\end{document}
